% steps/step10.tex
\section[گام دهم: اعتبارسنجی نهایی]{\faLock\ \ گام دهم: اعتبارسنجی نهایی}
\begin{stepbox}

\field{آیا مدل در شرایط واقعی عملکرد خوبی دارد؟}{
    بر اساس نتایج بنچمارک‌های مرجع مانند LiTS و عملکرد معماری‌های پیشرفته‌ای نظیر UNet++، مدل‌های یادگیری عمیق در شرایط کنترل‌شده آزمایشگاهی به دقت‌های بالایی (مانند \lr{Dice Score} بالای ۹۰٪ برای کبد و بین ۷۰ تا ۸۰٪ برای تومور) دست یافته‌اند. با این حال، انتقال مدل به محیط‌های بالینی واقعی نیازمند ارزیابی در شرایط عملیاتی است. به این منظور، سیستم در ابتدا به‌صورت موازی و در پس‌زمینه بیمارستان (\lr{Shadow Mode}) مستقر می‌شود تا بدون اثرگذاری مستقیم بر فرآیند درمان، خروجی‌های آن با تصمیمات واقعی رادیولوژیست‌ها مقایسه و صحه‌گذاری شود.
}

\field{ارزیابی روی داده‌های دیده‌نشده (External Validation)}{
    برای سنجش واقعی قدرت تعمیم‌پذیری، ارزیابی خارجی (\lr{External Validation}) با استفاده از داده‌های بیمارستان‌ها و اسکنرهایی که در فرآیند آموزش مدل حضور نداشته‌اند، الزامی است. این گام تضمین می‌کند که مدل نسبت به نویزهای خاص، ضخامت مقاطع و پروتکل‌های تصویربرداری یک مرکز درمانی دچار بیش‌برازش نشده و عملکرد پایدار خود را در مواجهه با توزیع‌های داده‌ای جدید حفظ می‌کند.
}

\field{بررسی خطاهای بحرانی و ریسک‌های بالقوه}{
    در سناریوهای تصویربرداری پزشکی، وزن خطاهای مختلف یکسان نیست. خطای منفی کاذب (\lr{False Negative}) که به معنای عدم شناسایی و نادیده گرفتن تومور توسط مدل است، تهدیدی حیاتی برای جان بیمار به شمار می‌رود و بسیار خطرناک‌تر از خطای مثبت کاذب (\lr{False Positive}) یا تشخیص اشتباه بافت سالم به‌عنوان تومور است. از این رو، ارزیابی نرخ منفی کاذب سیستم و مرزبندی‌های دقیق مرز ضایعات در سناریوهای دشوار بالینی، فاز کلیدی در تایید نهایی مدل خواهد بود.
}

\end{stepbox}
